\writeSectionFrame{\presentationTitle}{Das Problem}
\begin{frame}{Baumstrukturen}
	\begin{center}
 		\includegraphics[scale=0.25]{\BILDERPFAD/problem1.png}
 	\end{center}
\end{frame}

\note{
 	Wir haben ein Programm, das Emailadressen und Emailverteiler verwaltet. Dabei können Emailverteiler
 	weitere Emailverteiler enthalten, aber auch Emailadressen. Emailadressen können keine weiteren
 	Elemente mehr enthalten. 
}

\begin{frame}{Baumstrukturen}
	\begin{center}
 		\includegraphics[scale=0.25]{\BILDERPFAD/problem2.png}
 	\end{center}
\end{frame}

\note{
 	Also etwas allgemeiner ausgedrückt haben wir Komposita, die Elemente enthalten können, und Blätter,
 	bei denen der Baum endet. Wir entwickeln also eine Baumstruktur. Doch wie bildet man so etwas
 	einfach in einem Programm ab? Vor allem, wenn die Strukturen komplexer werden oder sogar während
 	der Laufzeit anpassbar sein sollen? 
}

\begin{frame}{Weitere Beispiele}
	\begin{itemize}
 		\item Dateien und (Unter)-Ordner
 		\item Objekte und (Unter)-Listen
 		\item Artikel und (Unter)-Kategorien
 		\item Mitarbeiter und (Unter)-Abteilungen
 		\item Projektergebnisse und (Teil)-Projekte
 		\item Komponenten und (Teil)-Komponenten
 		\item Rangordnung im Militär
 		\item Zusammengesetzte Grafiken
 	\end{itemize}
\end{frame}
